%==============================================================================
%  CAPÍTULO VI
%==============================================================================
\chapter{Procedimiento administrativo de renegociación de la Persona Deudora}
\label{cap:renegociacion}

\fichaCapitulo{Tramitación gratuita ante la \superir{}.}{Ingreso de boletas de
honorarios, audiencias de determinación, renegociación y ejecución.}

%==============================================================================
\bloqueTeorico
%==============================================================================

\subsection{La persona natural como sujeto de consumo financiero}
\pendiente{Desarrollar.}

\subsection{Gratuidad y desjudicialización como reductores del costo social}
\pendiente{Desarrollar: fundamento ético de la renegociación frente a la ejecución
forzosa.}

%==============================================================================
\bloqueNormativo
%==============================================================================

\subsection{Causales de admisibilidad}

El \art{260} \parencite{ley-20720} exige la concurrencia copulativa de los siguientes
requisitos:

\begin{enumerate}
  \item Dos o más obligaciones vencidas por más de \destacar{noventa días corridos},
  actualmente exigibles y \destacar{provenientes de obligaciones diversas}.
  \item Que el monto total sea superior a \uf{80}.
  \item Que la persona deudora no haya sido notificada de una demanda que solicite el
  inicio de un procedimiento concursal de liquidación simplificada, ni de cualquier otro
  juicio ejecutivo iniciado en su contra.
\end{enumerate}

\begin{foco}[La excepción laboral en la causal de inadmisibilidad]
El tercer requisito contiene una salvedad que suele pasarse por alto: la exclusión no
opera respecto de los juicios ejecutivos \destacar{de origen laboral}. Una persona
deudora notificada de una ejecución laboral conserva el acceso a la renegociación. La
verificación de causas por RUT debe, por tanto, distinguir la naturaleza de cada
ejecución y no limitarse a constatar su existencia.
\end{foco}

\subsection{Obligaciones excluidas del procedimiento}

El propio \art{260} sustrae determinadas obligaciones del ámbito de la renegociación:

\begin{enumerate}
  \item Los alimentos que se deben por ley, conforme al Título XVIII del Libro I del
  \cc{}, y la compensación económica del Párrafo 1.\textsuperscript{o} del Capítulo VII
  de la \leyn{19.947}{Ley de Matrimonio Civil}.
  \item Las obligaciones derivadas de delitos o cuasidelitos.
\end{enumerate}

\begin{alerta}[Advertencia obligada al cliente]
Estas exclusiones deben comunicarse en la primera reunión y consignarse por escrito. El
deudor que ingresa a la renegociación creyendo que reestructurará su deuda alimentaria
descubrirá lo contrario cuando ya no pueda revertirlo. Véase el capítulo
\ref{cap:entrevista}.
\end{alerta}

\subsection{Ampliación de la legitimación activa}

Incorporación de las personas naturales que emiten boletas de honorarios, por obra de la
\Lreforma{} \citeyearpar{ley-21563}. Sobre el impacto de la medida en el primer año de
vigencia, véase \textcite{superir-balance-primer-ano}.

\subsection{Efectos de la Resolución de Admisibilidad}

Desde la publicación de la Resolución de Admisibilidad y hasta el término del
procedimiento, el \art{264} produce los siguientes efectos:

\begin{enumerate}
  \item \textbf{Inmunidad ejecutiva} (\Nro 1). No puede solicitarse la liquidación
  forzosa ni voluntaria de la persona deudora, ni iniciarse en su contra juicios
  ejecutivos, ejecuciones de cualquier clase ni restituciones en juicios de arrendamiento.
  Para hacerla valer, el deudor acompaña copia autorizada de la resolución al tribunal
  respectivo, \destacar{y puede comparecer personalmente, sin patrocinio de abogado}. La
  ley precisa que sólo puede hacerse valer como excepción.

  \item \textbf{Suspensión de la prescripción} (\Nro 2). Se suspenden los plazos de
  prescripción extintiva de las obligaciones del deudor.

  \item \textbf{Cese del devengo de intereses moratorios} (\Nro 3). No continúan
  devengándose los intereses moratorios pactados en los actos y contratos vigentes.
\end{enumerate}

\begin{practica}[Un efecto de valor inmediato para el cliente]
El cese del devengo de intereses moratorios detiene el crecimiento del pasivo desde la
admisibilidad. En deudas de consumo con mora prolongada, este solo efecto puede
representar una porción sustancial del beneficio del procedimiento, y conviene explicarlo
al cliente en términos cuantitativos.
\end{practica}

\subsection{El acuerdo de ejecución y sus límites}

El acuerdo de ejecución puede contener un plan de reembolso, regulado en el \art{267},
sujeto a dos límites: no puede exceder de \destacar{seis meses} de duración, y la cuota
mensual no puede ser superior al \destacar{30\,\% de los ingresos del deudor}.

La \superir{} puede suspender la audiencia hasta por diez días.

\begin{alerta}[Corrección de una atribución errónea]
El plan de reembolso se regula en el \art{267} y no en el \art{264}, que trata de los
efectos de la resolución de admisibilidad. La confusión es frecuente en la literatura de
divulgación.
\end{alerta}

\subsection{Modificación del acuerdo de renegociación}

El acuerdo puede modificarse por una sola vez, dentro de los cinco años siguientes a la
resolución de admisibilidad.\verificar{Confirmar el precepto que regula la modificación y
sus requisitos.}

%==============================================================================
\bloquePractico
%==============================================================================

\subsection{Tramitación en la plataforma de la \superir{}}

\begin{flujograma}{Secuencia del procedimiento administrativo de renegociación}{cap06-iter}
  \node[hito] (ingreso) {INGRESO DE SOLICITUD EN EL PORTAL DE LA \superir{}};
  \node[nodo, below=of ingreso] (examen)
    {EXAMEN DE ADMISIBILIDAD POR EL FISCALIZADOR};
  \node[nodo, below=of examen] (resol)
    {RESOLUCIÓN DE ADMISIBILIDAD DE LA SOLICITUD};
  \node[nodo, below=of resol] (susp)
    {SUSPENSIÓN DE EJECUCIONES Y DE LA PRESCRIPCIÓN};
  \node[nodo, below=of susp] (pasivo)
    {AUDIENCIA DE DETERMINACIÓN DEL PASIVO\\[2pt]\normalfont Plazo de 15 días};
  \node[hito, below=of pasivo] (reneg)
    {AUDIENCIA DE RENEGOCIACIÓN};

  \coordinate (split) at ($(reneg.south)+(0,-6mm)$);
  \node[favorable, below left=12mm and 4mm of reneg.south] (aprueba)
    {APROBACIÓN DEL ACUERDO};
  \node[adverso, below right=12mm and 4mm of reneg.south] (rechaza)
    {RECHAZO DE LA PROPUESTA};
  \node[favorable, below=9mm of aprueba] (pub)
    {PUBLICACIÓN DEL ACUERDO EN EL \boletin{}};
  \node[adverso, below=9mm of rechaza] (ejec)
    {AUDIENCIA DE EJECUCIÓN};

  \draw[flecha] (ingreso) -- (examen);
  \draw[flecha] (examen) -- (resol);
  \draw[flecha] (resol) -- (susp);
  \draw[flecha] (susp) -- (pasivo);
  \draw[flecha] (pasivo) -- (reneg);
  \draw[flecha] (reneg.south) -- (split) -| (aprueba.north);
  \draw[flecha] (reneg.south) -- (split) -| (rechaza.north);
  \draw[flecha] (aprueba) -- (pub);
  \draw[flecha] (rechaza) -- (ejec);
\end{flujograma}

\verificar{Confirmar la duración de la audiencia de renegociación tras su ampliación por
la Ley 21.563 y el nombre vigente del portal de tramitación.}

\subsection{Preparación de las audiencias}
\pendiente{Desarrollar: documentación exigida, objeciones al pasivo, representación.}

%==============================================================================
\bloqueJurisprudencial
%==============================================================================

\subsection{Doctrina administrativa de la \superir{}}

Pautas de la \ncg{28} \citeyearpar{ncg-28} sobre acreditación de la existencia y monto de las obligaciones de
codeudores solidarios y avales dentro del proceso de renegociación.

\subsection{Recursos de protección por mantención de registros negativos}

Fallos que sancionan a instituciones financieras que mantienen vigentes registros
comerciales negativos respecto de deudores con acuerdos de renegociación firmes.
\verificar{Individualizar sentencias.}
